\documentclass[11pt]{article}
\usepackage{preamble}

\newcommand{\IconCheck}{[CHECK]}
\newcommand{\IconMic}{[MIC]}
\newcommand{\IconClipboard}{[CLIPBOARD]}

\newcommand{\phasetag}[1]{%
  {\small\itshape Tag: #1\par}
}

\newcommand{\phaseitems}[1]{%
  \begin{itemize}[leftmargin=1.5em,itemsep=2pt,topsep=2pt]
  #1
  \end{itemize}%
}

\newcommand{\teacheranswerbox}[2]{%
  \begin{callout}{#1}{calloutgreen}
  #2
  \end{callout}
}

\newcommand{\donowfullgraph}{%
\begin{tikzpicture}
\begin{axis}[
  width=0.43\linewidth,
  height=2.15in,
  xmin=-3,
  xmax=6,
  ymin=-1,
  ymax=8,
  axis lines=middle,
  grid=major,
  minor tick num=1,
  xlabel={$x$},
  ylabel={$y$},
  tick label style={font=\small},
  label style={font=\small}
]
\addplot[tealheader, line width=1.2pt, domain=-3:6, samples=2] {x+2};
\end{axis}
\end{tikzpicture}%
}

\newcommand{\donowholegraph}{%
\begin{tikzpicture}
\begin{axis}[
  width=0.43\linewidth,
  height=2.15in,
  xmin=-3,
  xmax=6,
  ymin=-1,
  ymax=8,
  axis lines=middle,
  grid=major,
  minor tick num=1,
  xlabel={$x$},
  ylabel={$y$},
  tick label style={font=\small},
  label style={font=\small}
]
\addplot[tealheader, line width=1.2pt, domain=-3:1.95, samples=2] {x+2};
\addplot[tealheader, line width=1.2pt, domain=2.05:6, samples=2] {x+2};
\addplot[
  only marks,
  mark=o,
  mark size=3pt,
  mark options={draw=tealheader, fill=white, line width=0.9pt}
] coordinates {(2,4)};
\end{axis}
\end{tikzpicture}%
}

\begin{document}

\packetheading{Lesson 4-3 · Period 1 · Teacher Edition}{Rational Expressions --- Equivalent \& Simplify}

\sectionbanner{OBJECTIVES}
{\small
\renewcommand{\arraystretch}{1.25}
\noindent\begin{tabular}{|p{1.8in}|p{4.8in}|}
\hline
\textbf{Math Objective} & Write equivalent rational expressions over a restricted domain, identify domain restrictions from factored denominators, and simplify rational expressions by canceling common factors. \\ \hline
\textbf{Language Objective} & Explain, in writing and orally using sentence frames, why a value is excluded from the domain and why we cancel factors but not terms. \\ \hline
\textbf{Essential Understanding} & A rational expression equals an equivalent simpler form ONLY when both share the same domain --- canceling a factor never erases the restriction it imposed. \\ \hline
\end{tabular}
}

\sectionbanner{TOPIC GOALS / ESSENTIAL QUESTION / MATERIALS}
{\small
\renewcommand{\arraystretch}{1.25}
\noindent\begin{tabular}{|p{1.8in}|p{4.8in}|}
\hline
\textbf{Topic Goals} & Write equivalent rational expressions, simplify, and identify domain restrictions from polynomial denominators. \\ \hline
\textbf{Essential Question} & Why does canceling a common factor not remove its domain restriction? \\ \hline
\textbf{Materials} & Student packet, pencil, highlighter. Teacher displays worked Example 1 then gradual-release into Example 2. \\ \hline
\end{tabular}
}

\begin{callout}{\IconPin\ PRE-FLIGHT NOTE (read before class)}{calloutpink}
Explore is 7 bank items: Try-It 1 \(\rightarrow\) \#20 \(\rightarrow\) \#22 \(\rightarrow\) \#24 \(\rightarrow\) Try-It 2 \(\rightarrow\) \#17 \(\rightarrow\) \#19 (extension). Paced for \(\sim\)4--5 min per item.\\
If pairs finish early, push them to the Optional Reinforcement: Model \& Discuss Parts B + C (verbal/TPS), then Practice \#25.\\
If pace slows, cut \#19 (extension) --- it's the DOK-2 stretch; the core DOK-1 arc (Try-It 1 \(\rightarrow\) \#20/\#22/\#24 \(\rightarrow\) Try-It 2 \(\rightarrow\) \#17) is the minimum viable set.
\end{callout}

\phasetag{notice \& wonder}
\frameworkphaseheader{Do Now}{2}{5}{%
\phaseitems{
  \item Project the graph (or use printed Do Now sheet).
  \item Circulate silently. Do NOT teach during this phase.
  \item Collect at 5 min sharp.
}%
}{%
\phaseitems{
  \item Write 1 notice + 1 wonder.
  \item Write one sentence about what the equation must be doing at the skipped \(x\)-value.
}%
}{%
\phaseitems{
  \item (None during the phase --- teacher harvests answers in Launch.)
}%
}{Solo teacher: circulate silently. Resist the urge to give hints.}

\bankitem{Do Now (4-3-savvas-model-discuss-lesson-4-3-launch)}{%
\textbf{EXPLORE \& REASON}

Consider the following graph of the function \(y=x+2\).

\begin{center}
\donowfullgraph
\end{center}

\begin{enumerate}[label=\textbf{\Alph*.}, leftmargin=1.6em,itemsep=3pt,topsep=3pt]
\item What is the domain of this function?
\item Sketch a function that resembles the graph, but restrict its domain to exclude 2.
\item Use Structure Consider the function you have sketched. What kind of function might have a graph like this? Explain.
\end{enumerate}

\textbf{Sample Student Work}

\begin{enumerate}[label=\textbf{\Alph*.}, leftmargin=1.6em,itemsep=3pt,topsep=3pt]
\item all real numbers.
\item
\begin{center}
\donowholegraph
\end{center}
\item An equation where the denominator equals zero when \(x=2\). For example, \(y=(2x-4)/(x-2)\).
\end{enumerate}
}{}

\begin{callout}{\IconCheck\ ANTICIPATED STUDENT RESPONSES}{calloutgreen}
NOTICE (typical): ``There's a hole in the graph.'' ``The line has a gap.'' ``One point is missing.''\\
WONDER (typical): ``Why is there a hole?'' ``What happens at that \(x\)?'' ``Can you have an equation that does this?''\\
SENTENCES: students might write ``the denominator equals zero at that \(x\)'' or ``the equation is undefined there.''\\
KEY MOVE in Launch: surface the insight that when the denominator \(= 0\), the expression is undefined --- this is a domain restriction.
\end{callout}

\phasetag{teacher models Example 1 + Example 2}
\frameworkphaseheader{Launch}{1--2}{12}{%
\phaseitems{
  \item Bridge from Do Now: ``You noticed the graph skips a point. That happens because the denominator equals zero at that \(x\)-value --- we call it a domain restriction.''
  \item Project Example 1 (Savvas Topic 4-3). Think aloud: factor numerator and denominator, identify what \(x\)-values make the denominator zero, write the equivalent expression with restriction stated.
  \item Project Example 2. Think aloud: factor completely, cancel the common factor, but keep the domain restriction from the canceled factor in the final answer.
  \item Pause \(1\times\) for 30-sec partner-check: ``What \(x\)-value is excluded in Example 1?''
  \item Do NOT solve Try-Its. Release promptly at minute 17.
}%
}{%
\phaseitems{
  \item Watch silently through Example 1.
  \item During partner-check: identify the excluded \(x\)-value.
  \item Copy the Equivalent + Simplify Rules card into their page margin.
  \item Note the COMMON ERROR warning: cancel only factors, never terms.
}%
}{%
\phaseitems{
  \item ``What values make the denominator zero?''
  \item ``Can I cancel this factor? Is it a factor of the entire numerator?''
  \item ``After canceling, does the domain restriction disappear?''
}%
}{Project, think aloud. Do NOT give students Try-Its to work until minute 17.}

\begin{callout}{\IconMic\ TEACHER SCRIPT}{calloutblue}
1. ``From the Do Now, you noticed the graph skips a point. That \(x\)-value is excluded from the domain because the denominator equals zero there.''\\
2. Example 1 walk-through: ``Watch me. I factor top and bottom. Any \(x\) that makes the denominator zero is restricted --- I write \(x \ne\) that value next to my answer.''\\
3. Example 2 walk-through: ``Now I simplify. I factor, find the matching factor in numerator and denominator, and cancel it. But --- critical --- the restriction from that canceled factor STAYS. Canceling doesn't erase the hole.''\\
4. Common-error flag: ``Don't cancel terms that are added or subtracted. Factor FIRST, then cancel only matching factor pairs.''\\
5. Release: ``Now you try. Try-It 1 asks for an equivalent expression. Try-It 2 asks you to simplify completely.''
\end{callout}

\phasetag{Think-Pair-Share + Practice}
\frameworkphaseheader{Explore}{1--2}{33}{%
\phaseitems{
  \item 5 laps circulation across 33 min (\(\sim\)6--7 min each).
  \item Lap 1: Try-It 1 --- press pairs to factor before writing the equivalent form.
  \item Lap 2: \#20 + \#22 --- check that pairs identify domain restrictions correctly.
  \item Lap 3: \#24 + Try-It 2 --- press pairs to factor completely before canceling.
  \item Lap 4: \#17 --- verify pairs state the domain restriction alongside the simplified form.
  \item Lap 5: \#19 (extension) --- stretch. Cut if pace slows.
  \item Do not give answers. Use sentence frame to press justification.
}%
}{%
\phaseitems{
  \item Pairs move in order through 7 items: 2 Try-Its + 5 Practice.
  \item For each item: factor completely FIRST, then cancel, then state restriction.
  \item For Practice \#19 (simplify with domain restrictions): show all steps.
  \item Justify with sentence frame before writing.
}%
}{%
\phaseitems{
  \item Did you factor completely before canceling?
  \item What \(x\)-values make the denominator zero?
  \item After simplifying, what is the domain restriction?
  \item Can you use the sentence frame to explain your answer?
  \item For \#19: what is the domain restriction from the canceled factor?
}%
}{Solo teacher: 5 laps. Press pairs to justify using the rule card. No answers.}

\bankitem{Try It 1 (4-3-savvas-try-it-1-lesson-4-3)}{%
Write an expression equivalent to \((x-4)/(x)\) over the domain \(\{x\mid x\ne 0 \text{ or } -2\}\).
}{}

\teacheranswerbox{\IconCheck\ ANSWER / ANTICIPATED TALK}{%
\textbf{Answer key (from source):} \((x^2-2x-8)/(x^2+2x)\)
}

\bankitem{Practice \#20 --- Identify equivalent (4-3-savvas-q20)}{%
Write an equivalent expression over the given domain.

\[
\frac{x(x-2)}{x-5} \text{ for all } x \text{ except } 5 \text{ and } -6
\]
}{}

\teacheranswerbox{\IconCheck\ ANSWER / ANTICIPATED TALK}{%
\textbf{Answer key (from source):} \((x^3+4x^2-12x)/(x^2+x-30)\)
}

\bankitem{Practice \#22 (4-3-savvas-q22)}{%
What is the simplified form of the rational expression? What is the domain?

\[
\frac{y^2-5y-24}{y^2+3y}
\]
}{}

\teacheranswerbox{\IconCheck\ ANSWER / ANTICIPATED TALK}{%
\textbf{Answer key (from source):} \((y-8)/(y)\) for all \(y\) except \(0\) and \(-3\).
}

\bankitem{Practice \#24 (4-3-savvas-q24)}{%
What is the simplified form of the rational expression? What is the domain?

\[
\frac{x^2+8x+15}{x^2-x-12}
\]
}{}

\teacheranswerbox{\IconCheck\ ANSWER / ANTICIPATED TALK}{%
\textbf{Answer key (from source):} \((x+5)/(x-4)\), \(x\ne -3,4\).
}

\bankitem{Try It 2 (4-3-savvas-try-it-2-lesson-4-3)}{%
Simplify each expression and state the domain.

\begin{enumerate}[label=\textbf{\alph*.}, leftmargin=1.6em,itemsep=3pt,topsep=3pt]
\item \(\dfrac{x^2+2x+1}{x^3-2x^2-3x}\)
\item \(\dfrac{x^3+4x^2-x-4}{x^2+3x-4}\)
\end{enumerate}
}{}

\teacheranswerbox{\IconCheck\ ANSWER / ANTICIPATED TALK}{%
\textbf{Answer key (from source):} a. \((x+1)/(x(x-3))\) for all \(x\) except \(0\), \(-1\), and \(3\). b. \(x+1\) for all \(x\) except \(1\) and \(-4\).
}

\bankitem{Practice \#17 --- Simplify (4-3-savvas-q17)}{%
\textbf{Construct Arguments} Explain how you can tell whether a rational expression is in simplest form.
}{}

\teacheranswerbox{\IconCheck\ ANSWER / ANTICIPATED TALK}{%
\textbf{Answer key (from source):} A rational expression is in simplest form when the numerator and denominator have no common factors other than 1.
}

\bankitem{Practice \#19 --- Simplify with domain restrictions (4-3-savvas-q19)}{%
\textbf{Reason} If the denominator of a rational expression is \(x^3+3x^2-10x\), what value(s) must be restricted from the domain for \(x\)?
}{}

\teacheranswerbox{\IconCheck\ ANSWER / ANTICIPATED TALK}{%
\textbf{Answer key (from source):} \(-5\), \(0\), \(2\)
}

\phasetag{academic conversation}
\frameworkphaseheader{Share / Summary}{2}{5}{%
\phaseitems{
  \item Cold-call 2 pairs to read their sentence-frame sentence.
  \item Write one student-generated sentence on board.
  \item Preview P2: ``Next class we'll multiply and divide rational expressions --- the same factor-and-cancel logic applies.''
}%
}{%
\phaseitems{
  \item Presenting pair reads their sentence.
  \item Class signals thumbs up/sideways.
  \item Note: next class extends this to multiplication and division.
}%
}{%
\phaseitems{
  \item ``What's a sentence that captures today's rule about domain restrictions?''
  \item ``Why can't we just ignore the restriction after we cancel?''
}%
}{Cold-call 2 pairs. Keep it brief --- 5 min.}

\phasetag{summary recap (not graded)}
\frameworkphaseheader{Exit Ticket}{1--2}{2}{%
\phaseitems{
  \item Collect at door.
  \item Scan the fill-in items --- misconceptions about domain restrictions tell you what to re-hit in P2 Do Now.
}%
}{%
\phaseitems{
  \item Complete the fill-in stems.
  \item Be honest about confusion.
}%
}{%
\phaseitems{
  \item Did students correctly identify domain restrictions from the original denominator?
}%
}{Collect. No grade.}

\begin{callout}{\IconClipboard\ EXIT TICKET --- Student Form}{calloutyellow}
To simplify \((x^2+3x-10)/(x^2-4)\) I first factor to \_\_\_.\\
The restricted domain values are \(x \ne\) \_\_\_.\\
After canceling the factor \_\_\_, my simplified expression is \_\_\_.\\
One biggest thing I learned today: \_\_\_\_\_\_\_\_\_\_\_\_\_\_\_\_\_\_\_\_\_\_\_\_\_\_
\end{callout}

\clearpage

\begin{callout}{\IconClock\ PERIOD A \& F --- 65-MIN CLASS NOTES}{calloutpink}
Both sections get 65 min for P1. Use the extra 10 min on Explore, NOT Launch.\\
Suggested add: project Savvas Practice \#22 and \#24 on screen for 2 more TPS rounds if pairs finish fast.\\
Or: extend Model \& Discuss Parts B + C with ``write one more rational expression that has the same restriction --- then simplify it.''
\end{callout}

\begin{callout}{\IconPuzzle\ MODIFICATIONS / SUPPORT FOR IEP STUDENTS}{calloutiep}
\textbullet\ Provide the Equivalent + Simplify Rules card as a sticker/cut-out on their desk.\\
\textbullet\ For Try-It 1, allow students to factor with a reference list of factoring patterns.\\
\textbullet\ Accept verbal sentence-frame completions in lieu of written.
\end{callout}

\begin{callout}{\IconGlobe\ SUPPORTS FOR ELL STUDENTS}{calloutell}
\textbullet\ Sentence frames printed on student packet (don't skip).\\
\textbullet\ Word card: ``domain restriction'' = an \(x\)-value that makes the denominator zero; excluded from the domain.\\
\textbullet\ ``cancel'' = divide out a common factor from numerator and denominator.\\
\textbullet\ ``factor'' = rewrite as a product of simpler expressions.\\
\textbullet\ Pair ELL students with English-dominant partners for TPS.
\end{callout}

\end{document}
